%% capitulos/2-fundamentacao-teorica.tex

\chapter{Fundamentação Teórica}
\label{cap:fundamentacao}

Este capítulo apresenta os conceitos teóricos que fundamentam
esta dissertação.

\section{Tema Principal}
\label{sec:tema-principal}

Apresente os conceitos centrais da sua pesquisa, com embasamento
bibliográfico adequado \cite{exemplo2023}.

% Exemplo de citação direta curta (≤ 3 linhas, entre aspas):
Segundo \citeonline{exemplo2023}, ``texto citado diretamente entre
aspas, com até três linhas, inserido no corpo do parágrafo''.

% Exemplo de citação direta longa (> 3 linhas):
% Use o ambiente 'citacao' do abntex2:
\begin{citacao}
Texto com mais de três linhas que deve ser recuado 4 cm da margem
esquerda, com fonte tamanho 10, espaço simples, sem aspas.
O recuo e formatação são automáticos com este ambiente.
\cite[p. 15]{exemplo2023}
\end{citacao}

\section{Subtema 1}
\label{sec:subtema1}

Texto da seção \ldots

\subsection{Subseção 1.1}
\label{subsec:subsecao11}

Texto da subseção \ldots

% Exemplo de figura:
\begin{figure}[htb]
  \caption{Título descritivo da figura}
  \label{fig:exemplo}
  \begin{center}
    \includegraphics[width=0.6\textwidth]{figuras/exemplo}
  \end{center}
  \legend{Fonte: Autor (2025)}
\end{figure}

A Figura~\ref{fig:exemplo} ilustra \ldots

% Exemplo de tabela:
\begin{table}[htb]
  \caption{Título descritivo da tabela}
  \label{tab:exemplo}
  \begin{center}
    \begin{tabular}{lcc}
      \toprule
      \textbf{Coluna 1} & \textbf{Coluna 2} & \textbf{Coluna 3} \\
      \midrule
      Dado A & 10 & 20 \\
      Dado B & 15 & 30 \\
      \bottomrule
    \end{tabular}
  \end{center}
  \legend{Fonte: Elaborada pelo autor (2025)}
\end{table}

A Tabela~\ref{tab:exemplo} apresenta \ldots
